\documentclass[11pt,letterpaper]{article}
\usepackage[letterpaper, left=0.4in, right=0.4in, top=0.4in, bottom=0.4in]{geometry}
\usepackage{mathptmx}
\usepackage{enumitem}
\usepackage[hidelinks]{hyperref}
\usepackage{titlesec}
\pagestyle{empty}
\titleformat{\section}{\large\bfseries\uppercase}{}{0em}{}[\titlerule]
\titlespacing{\section}{0pt}{3pt}{2pt}
\setlist[itemize]{leftmargin=*, itemsep=0.3pt, topsep=1pt, parsep=0pt, partopsep=0pt}
\begin{document}
\begin{center}
{\LARGE \textbf{SHRENIK PILLAI}} \\ \vspace{2pt}
\small US CITIZEN \quad Richardson, TX \\ \vspace{1pt}
\small (469) 545-8106 \quad \href{mailto:rshrenik811@gmail.com}{rshrenik811@gmail.com} \quad \href{https://linkedin.com/in/shrenik-r}{linkedin.com/in/shrenik-r} \quad \href{https://github.com/Shrenik811rp}{github.com/Shrenik811rp}
\end{center}
\section*{SUMMARY}
\noindent\small Software engineer with 3 years of experience architecting distributed systems on AWS and Linux environments. Built infrastructure automation pipelines using Python, Ansible, and Jenkins, reducing manual operations and accelerating deployment cycles. Containerized microservices with Docker and managed Kubernetes clusters for production workloads. Experience optimizing system telemetry and monitoring SLIs to reduce incident resolution time. Skilled in translating infrastructure requirements into scalable cloud solutions across compute, storage, and networking domains.
\section*{WORK EXPERIENCE}
\noindent\small \textbf{o9 Solutions, Inc.} \hfill Jan 2023 to Aug 2025 \\
\textit{DevOps & Infrastructure Systems Engineer} \hfill \textit{Remote}
\begin{itemize}
\item Built observability dashboards using Python and SQL to monitor distributed system SLIs (latency, error rates), enabling proactive anomaly detection that reduced Mean Time to Resolve (MTTR) by 14% across production environments.
\item Designed and automated CI/CD deployment workflows using Python and Ansible, reducing manual toil by 10% and cutting Jenkins pipeline lead time for application releases.
\item Containerized microservices with Docker and orchestrated deployments across Kubernetes clusters, standardizing configuration management through Git to reduce environment drift by 12%.
\item Architected Linux-based system administration automation for log rotation, resource monitoring, and backup processes across 50+ virtualized nodes, improving system availability.
\item Tracked SLA compliance metrics and reported weekly platform reliability status to cross-functional stakeholders, coordinating CI/CD upgrades to maintain project delivery timelines.
\end{itemize}
\noindent\small \textbf{mokSa.ai} \hfill Jan 2022 to Apr 2022 \\
\textit{Computer Vision & Infrastructure Engineer} \hfill \textit{Remote}
\begin{itemize}
\item Developed Python data processing pipelines for 50GB+ of daily sensor telemetry, implementing CUDA-accelerated preprocessing with NumPy and OpenCV for real-time computer vision workloads.
\item Wrote validation scripts in Python and Shell to ensure system quality compliance for deep learning inference pipelines, improving QA tracking accuracy by 8%.
\item Analyzed GPU and CPU resource utilization in Linux environments to identify performance bottlenecks, optimizing compute efficiency through better job scheduling and resource allocation.
\end{itemize}
\section*{PROJECT EXPERIENCE}
\noindent\small \textbf{Automated MLOps & Evaluation Pipeline} \hfill Jan 2026 to May 2026
\begin{itemize}
\item Built telemetry analytics pipelines in Python and Shell to monitor distributed training jobs, tracking real-time KPI anomalies and preventing execution drift by 10%.
\item Constructed CI/CD testing harnesses using Git and Jenkins, standardizing pre-merge validation with containerized test environments in Docker to improve developer throughput by 12%.
\end{itemize}
\section*{LEADERSHIP EXPERIENCE}
\noindent\small \textbf{Google Developer Student Club} \hfill Jan 2022 to May 2023 \\
\textit{Technical Lead} \hfill \textit{Academic Chapter}
\begin{itemize}
\item Led technical workshops on Python, Git, and Linux for a 20+ member engineering team, standardizing development practices and tracking project milestones.
\item Troubleshot environment bottlenecks and system conflicts on Linux platforms, ensuring project deliverables remained on schedule.
\end{itemize}
\section*{EDUCATION}
\noindent\small \textbf{The University of Texas at Dallas}, Richardson, TX \hfill Aug 2025 to May 2027 \\
\textit{Master of Science in Business Analytics and Artificial Intelligence} \hfill \textit{GPA: 3.50/4.00} \\
\noindent\small \textbf{Amrita Vishwa Vidyapeetham} \hfill Aug 2019 to Oct 2023 \\
\textit{B.S. in Electronics & Computer Engineering} \hfill \textit{GPA: 3.34/4.00}
\section*{TECHNICAL SKILLS}
\noindent\small \textbf{Cloud & Infrastructure:} AWS, Azure, Kubernetes, Docker, Linux Administration, Ansible, Terraform \\
\textbf{Programming Languages:} Python, C++, Java, Go, SQL, Shell Scripting \\
\textbf{CI/CD & Automation:} Jenkins, GitLab, Git, Infrastructure-as-Code \\
\textbf{ML/DL Frameworks:} NumPy, SciPy, TensorFlow, cuDNN, CUDA, SLURM \\
\textbf{Observability:} System Telemetry, SLA/SLI Metrics, Incident Response, MTTR Reduction
\section*{CERTIFICATIONS}
\noindent\small \textbf{Microsoft Certified:} Azure Fundamentals (AZ-900) \\
\textbf{DeepLearning.AI:} Introduction to TensorFlow for AI, ML, and Deep Learning
\end{document}